\setcounter{chapter}{5}
\chapter{Testing and Verification}
\label{ch:testing}

\section{Testing Strategy}
\label{sec:test-strategy}

Software testing for SmartHire was conducted across five levels: unit testing, integration testing, system testing, performance testing, and security testing. Each level targets a different scope of the application and uses different techniques to verify correctness.

\textbf{Unit testing} verifies individual functions and components in isolation. Service-layer functions (e.g., score computation, input validation) were tested with known inputs and expected outputs.

\textbf{Integration testing} verifies that multiple modules work correctly together. For example, testing the application submission flow requires the authentication module, resume upload service, ATS scoring service, and database layer to cooperate correctly.

\textbf{System testing} evaluates the complete application from the user's perspective, simulating end-to-end workflows such as ``candidate applies for a job, completes all assessments, and receives an offer.''

\textbf{Performance testing} measures response times under varying load conditions to identify bottlenecks.

\textbf{Security testing} attempts to exploit common web vulnerabilities (XSS, CSRF, SQL injection, unauthorised data access) to verify that the application's defences are effective.

\section{Unit Test Cases}
\label{sec:unit-tests}

\begin{table}[htbp]
\centering
\caption{Unit test cases}
\label{tab:unit-tests}
\small
\begin{tabular}{p{0.6in}p{1.4in}p{1.5in}p{1.5in}p{0.5in}}
\toprule
\textbf{ID} & \textbf{Component} & \textbf{Test Description} & \textbf{Expected Result} & \textbf{Status} \\
\midrule
UT-01 & Auth service & Valid email/password login & JWT token returned & Pass \\
UT-02 & Auth service & Invalid password & Error 401 returned & Pass \\
UT-03 & Auth service & Empty email field & Validation error & Pass \\
UT-04 & Job service & Create job with valid data & Job record created with draft status & Pass \\
UT-05 & Job service & Create job missing title & Validation error & Pass \\
UT-06 & ATS scoring & Resume with matching skills & Score $\geq$ 7 & Pass \\
UT-07 & ATS scoring & Resume with no matching skills & Score $\leq$ 3 & Pass \\
UT-08 & MCQ scoring & All correct answers & Score = 100\% & Pass \\
UT-09 & MCQ scoring & No correct answers & Score = 0\% & Pass \\
UT-10 & MCQ scoring & Partial correct answers & Score proportional & Pass \\
UT-11 & Coding scorer & All test cases pass & Score = 100\% & Pass \\
UT-12 & Coding scorer & Compilation error & Score = 0\%, error message & Pass \\
UT-13 & Offer generator & Valid input data & PDF generated without errors & Pass \\
UT-14 & Offer generator & Missing salary field & Validation error & Pass \\
\bottomrule
\end{tabular}
\end{table}

\section{Integration Test Cases}
\label{sec:integration-tests}

\begin{table}[htbp]
\centering
\caption{Integration test cases}
\label{tab:integration-tests}
\small
\begin{tabular}{p{0.6in}p{1.8in}p{1.5in}p{1.5in}p{0.5in}}
\toprule
\textbf{ID} & \textbf{Workflow} & \textbf{Test Description} & \textbf{Expected Result} & \textbf{Status} \\
\midrule
IT-01 & Registration $\to$ Login & Register then login & Redirect to portal & Pass \\
IT-02 & Job Create $\to$ Publish & Create draft then publish & Job appears in directory & Pass \\
IT-03 & Apply $\to$ ATS Score & Submit application & ATS score computed and stored & Pass \\
IT-04 & MCQ Start $\to$ Submit & Begin and complete exam & Score recorded in application & Pass \\
IT-05 & Coding $\to$ Score & Submit code solution & Test results and score stored & Pass \\
IT-06 & Interview $\to$ Scorecard & Record and submit & AI scorecard generated & Pass \\
IT-07 & Pipeline $\to$ Reject & Move card to rejected & Rejection stage recorded & Pass \\
IT-08 & Offer $\to$ Accept & Generate and accept offer & Application status = hired & Pass \\
\bottomrule
\end{tabular}
\end{table}

\section{System Test Cases}
\label{sec:system-tests}

\begin{table}[htbp]
\centering
\caption{System-level end-to-end test cases}
\label{tab:system-tests}
\small
\begin{tabular}{p{0.6in}p{2in}p{1.3in}p{1.3in}p{0.5in}}
\toprule
\textbf{ID} & \textbf{Scenario} & \textbf{Expected Outcome} & \textbf{Actual Outcome} & \textbf{Status} \\
\midrule
ST-01 & Complete candidate journey (apply through all stages to offer) & Offer received & Offer received & Pass \\
ST-02 & Recruiter creates job, reviews applicants, hires one & Applicant hired & Applicant hired & Pass \\
ST-03 & Multi-tenant isolation: Recruiter A cannot see Recruiter B's candidates & No cross-tenant data & No cross-tenant data & Pass \\
ST-04 & Timer expiry during MCQ exam & Auto-submission occurs & Auto-submission occurs & Pass \\
ST-05 & Candidate rejected at coding stage views feedback & Rejection stage shown & Rejection stage shown & Pass \\
ST-06 & Admin dashboard shows platform-wide metrics & Aggregated data displayed & Aggregated data displayed & Pass \\
\bottomrule
\end{tabular}
\end{table}

\section{Performance Testing}
\label{sec:perf-tests}

Performance was measured for key operations under normal load conditions (single concurrent user) and moderate load (ten concurrent simulated users).

\begin{table}[htbp]
\centering
\caption{Performance test results}
\label{tab:perf-tests}
\small
\begin{tabular}{p{2.2in}ccc}
\toprule
\textbf{Operation} & \textbf{Single User} & \textbf{10 Users} & \textbf{Target} \\
\midrule
Page load (dashboard)          & 0.8 s  & 1.2 s  & $<$ 2.0 s \\
Job listing query              & 120 ms & 180 ms & $<$ 500 ms \\
Resume upload + ATS scoring    & 2.1 s  & 3.4 s  & $<$ 5.0 s \\
MCQ submission + scoring       & 95 ms  & 140 ms & $<$ 500 ms \\
Code execution (simple script) & 1.8 s  & 2.9 s  & $<$ 5.0 s \\
AI interview evaluation        & 4.2 s  & 6.1 s  & $<$ 10.0 s \\
Offer PDF generation           & 1.1 s  & 1.3 s  & $<$ 3.0 s \\
\bottomrule
\end{tabular}
\end{table}

All operations met their performance targets. The AI interview evaluation took the longest due to the round-trip to the Gemini API, but this is an asynchronous operation that does not block the candidate's interaction.

\section{Security Testing}
\label{sec:security-tests}

\begin{table}[htbp]
\centering
\caption{Security test results}
\label{tab:security-tests}
\small
\begin{tabular}{p{2in}p{2in}p{1.5in}}
\toprule
\textbf{Vulnerability} & \textbf{Test Method} & \textbf{Result} \\
\midrule
Cross-Site Scripting (XSS) & Injected script tags in form fields & Sanitised; no execution \\
SQL Injection & Injected SQL in search parameters & Parameterised queries; no effect \\
Unauthorised API access & Called API without valid JWT & 401 Unauthorized returned \\
Cross-tenant data access & Queried applications with another recruiter's ID & Empty result set \\
Direct object reference & Accessed application by ID belonging to another candidate & 403 Forbidden returned \\
\bottomrule
\end{tabular}
\end{table}

\section{Bug Report Summary}
\label{sec:bugs}

During testing, three notable bugs were identified and resolved:

\begin{enumerate}[leftmargin=1.5em]
  \item \textbf{Timer drift in MCQ engine}: The client-side countdown timer drifted by up to two seconds over a thirty-minute exam due to JavaScript's non-deterministic \texttt{setInterval} timing. This was resolved by computing remaining time from the difference between the current timestamp and the exam start timestamp rather than decrementing a counter.

  \item \textbf{Supabase client re-instantiation}: Rapid navigation between pages caused multiple Supabase client instances to be created, leading to authentication state inconsistencies. The singleton pattern described in Chapter~5 resolved this issue.

  \item \textbf{Kanban drag-and-drop race condition}: Rapidly dragging a candidate card between columns could trigger multiple simultaneous API calls, potentially resulting in inconsistent stage values. A debounce mechanism was added to prevent duplicate submissions.
\end{enumerate}
